The Jefferson method occupies a unique position in the history of
apportionment: it is simultaneously the first U.S.\ congressional
apportionment formula and the world's most widely used electoral
seat allocation method.
Its floor rounding rule --- simple, deterministic, and producing
no fractional-seat complications --- has an intuitive appeal that
explains both Jefferson's original advocacy and d'Hondt's independent
rediscovery.

The formal equivalence between Jefferson and d'Hondt (Theorem~\ref{thm:dhondt-equiv})
has important implications for the comparative electoral systems literature:
the academic debates about d'Hondt's large-party bias in PR systems
are mathematically equivalent to the historical U.S.\ debates about
Jefferson's large-state bias in apportionment.
The resolution adopted in both contexts was the same:
move toward arithmetic-mean rounding (Webster/Sainte-Lagu\"e),
which is neutral, or toward ceiling rounding (Adams), which
protects small entities.

For the \textsc{bisect} platform, the Jefferson/d'Hondt implementation
serves two purposes.
First, it enables counterfactual analysis of U.S.\ apportionment:
what would the 2020 congressional map look like if Jefferson's method
had never been replaced?
Second, it enables comparative PR analysis: how do different seat
allocation formulas translate vote distributions to seat allocations
for countries using d'Hondt?

The large-state-favoring bias of Jefferson is documented both theoretically
(Theorem~\ref{thm:jefferson-bias}) and historically (Section~\ref{sec:history}).
The 1832 controversy that led to Jefferson's abandonment was a
genuine mathematical dispute about quota compliance, not merely
a political conflict.
Webster's argument that floor rounding could give a state more than
its upper quota was mathematically correct, and Congress ultimately
agreed.
Understanding this history provides essential context for evaluating
modern apportionment reform proposals and for legal analysis of
apportionment challenges under Article~I.
